\section{Perspectives et tendances}\label{SEC:Perspectives et tendances}

\subsection{Miniaturisation et intégration RF}

Technologies CMOS, SiGe, GaAs.

Modules radar-on-chip (ex. TI AWR2944, Infineon BGT60 pour FMCW-SAR de courte portée).


\subsection{Intégration matérielle et logicielle}

Co-design matériel/logiciel pour traitement SAR temps réel.

Architectures SoC et FPGA reconfigurables (Xilinx Zynq, Intel SoC FPGA).


\subsection{Intelligence embarquée}

SAR couplé à l’IA (détection de cibles, compression intelligente).

Fusion de capteurs (SAR + EO/IR + Lidar).


\subsection{Limites technologiques actuelles}


Stabilité de phase, volume de données, puissance embarquée.


\subsection{Perspectives}


SAR à ouverture synthétique multi-bande, SAR 3D/Interférométrique.

Architectures reconfigurables, radar cognitif, IA embarquée.

Tendance vers le tout intégré (RF + numérique sur une seule puce).
